\section{Conclusion}
The proposed full digital workflow develops rasters from printhead and fluid inputs, a target font and production requirements. It links structure discovery, interacting-flight optimization, shared dependency fitting, integer-command replay and controller export, followed by printer qualification. Reusable device fields and compiled particle operations make repeated candidate evaluation practical, while controller-table coverage provides an explicit development target.

Two application cases demonstrate this workflow at the reported operating points. Fixed-command ablations quantify the effect of omitted interactions, Euler--RK4 comparisons assess numerical sensitivity, and matched CPU/GPU measurements characterize execution cost. The experimental comparison is against released ten-print coordinate means; its scope and the larger high-speed residuals are reported directly.

The principal result is a computational raster-development procedure with realizable outputs and traceable assessment. Broader device transfer and industrial time or cost savings remain separate evaluation questions.

The deposited-coordinate comparison gives deflection RMSE of 16.1--21.6 micrometers and substrate-direction RMSE of 21.0--31.2 micrometers against the experimental ten-print means. The high-speed case has the larger residuals, which remain an accuracy limit of the present model and operating-point assessment.

\section*{Data and code availability}
The software is distributed under the Apache License 2.0. The accompanying repository contains the simulation implementation, optimization workflows, supplied model inputs, released experimental coordinate means, numerical evidence and manuscript sources. Raw laboratory records and replicate-level experimental data remain confidential. The original CFD mesh and complete field-generation record are outside the supplied material.
